En este trabajo se estudiaron distintos métodos de meta-learning y domain adaptation aplicados a tareas de clasificación de dígitos bajo escenarios de cambio de dominio con pocos ejemplos.

En primer lugar, la evaluación de ProtoNet, MAML y ProtoMAML sobre MNIST-M y SVHN permitió observar con claridad el impacto del cambio de dominio. Aunque MNIST-M mantiene cierta similitud estructural con MNIST, sobre este dominio se reduce el desempeño. En SVHN, esta caída es aún más marcada, ya que las imágenes provienen de un contexto con mayor variabilidad de distintos aspectos. 

Entre los métodos de meta-learning comparados, ProtoNet y ProtoMAML obtuvieron mejores resultados que MAML en los dominios externos. Este comportamiento sugiere que los enfoques basados en representaciones prototípicas son más robustos frente al cambio de dominio que una adaptación basada únicamente en gradientes con muy pocos ejemplos. Además, ProtoMAML logró resultados similares o superiores a ProtoNet aun habiendo sido entrenado en la mitad de iteraciones y con menor cantidad de ejemplos por clase, lo que indica que la combinación entre inicialización prototípica y adaptación interna puede ser una estrategia efectiva para escenarios few-shot.

Por otro lado, los experimentos de domain adaptation mostraron que DANN mejora de manera significativa la transferencia de MNIST a MNIST-M respecto del baseline. Mientras que el baseline alcanza una accuracy alta sobre MNIST pero cae considerablemente en MNIST-M, DANN mantiene un rendimiento similar en el dominio fuente y obtiene una mejora importante en el dominio objetivo ($30\%$). Este resultado confirma que el entrenamiento adversarial mediante la Gradient Reversal Layer favorece la construcción de embeddings más invariantes al dominio, permitiendo que el modelo aproveche información no etiquetada del dominio target sin necesidad de realizar un ajuste supervisado completo.

El análisis de los embeddings mediante t-SNE acompaña esta interpretación: los modelos entrenados únicamente sobre MNIST tienden a organizar mejor las muestras del dominio fuente, pero presentan mayor dispersión y mezcla cuando se incorporan muestras de MNIST-M o SVHN. En cambio, DANN logra una mejor alineación entre MNIST y MNIST-M, lo que explica su mejora en accuracy sobre el dominio objetivo. De todas formas, el baseline con fine-tuning obtiene la separación más clara, algo esperable porque utiliza etiquetas del dominio target y funciona como una cota superior supervisada.

Como trabajo futuro, sería interesante incorporar técnicas de data augmentation para aumentar la variabilidad visual durante el entrenamiento y mejorar la robustez de los embeddings frente a cambios de dominio. También podrían evaluarse estrategias de adaptación más avanzadas, diferentes arquitecturas de encoder o combinaciones entre meta-learning y domain adaptation, con el objetivo de obtener modelos que no solo aprendan con pocos ejemplos, sino que además puedan generalizar mejor a dominios visualmente distintos.